\subsection{Equal CPU time on the standard benchmark}
\label{sec:eqcpu}
\paragraph{Design} The registered comparison of Section~\ref{sec:fanjul} is at equal iterations. We repeated it at equal CPU time in B30: the cost ratio $K$ of $1000$ iterations of each method to $1000$ iterations of random choice was measured in sequence on two tuning instances of each of the $28$ classes, random choice received $\mathrm{round}(1000K)$ iterations, and the validated action received $\mathrm{round}(1000K/K_{\mathrm{val}})$ iterations, so that it competes with each learner on the same CPU time. The median $K$ is $0.74$ for the standard Q-Learning, between $0.98$ and $1.07$ for UCB1, the policy gradient, the ALNS roulette, the optimistic Q-Learning and the FRRMAB-style design, $1.48$ for the episodic design and $1.97$ for the validated action (a local-search action in all classes), which therefore runs about half as many iterations as a learner at equal CPU time. Seven predictions $E_1$--$E_7$ and the hashes were recorded before the test run; two held (the others are discussed below).

\paragraph{Results} The cost effect of the standard Q-Learning is real. At equal iterations it loses $0.46$ points to random choice; at equal CPU time it is level with it ($-0.03$), because random choice receives only about three quarters of the iterations: its mean cost effect is $+0.44$ points, and it beats random choice beyond the margin in two of the seven types ($U[1,100]$ and $U[10,100]$: $+0.55$ and $+0.42$ points, with decision effects of $-0.57$ and $-0.42$), the types where the cheapest action, the one it repeats, happens to be good ($E_4$ failed: at most one type was predicted). The other learners have cost effects of at most $0.11$ points in absolute value, except the episodic design ($-0.55$: it costs $1.48$ per iteration) ($E_2$ failed on that learner). The validated action, with half the iterations, still beats the learners in $44$ of the $49$ learner--type pairs (confirmed; $E_5$ asked for $45$) and by $0.88$ points on average over the learner that comes closest in each type ($E_7$ asked for $1.0$). Two pairs go the other way, which we had predicted not to happen ($E_6$): on correlated machines the ALNS roulette and the FRRMAB-style design beat the validated action at equal CPU time, by $0.29$ and $0.53$ points. The statement that a validated action is never beaten beyond the margin by an online learner therefore holds at equal iterations on every host and, at equal CPU time on the standard benchmark, in $47$ of $49$ pairs, the exceptions being correlated machines, where the local-search action is the most expensive (median $K=2.50$, against $1.62$--$2.21$ in the other types).
